% Part: lambda-calculus
% Chapter: introduction
% Section: primitive-recursion

\documentclass[../../../include/open-logic-section]{subfiles}

\begin{document}

\olfileid{lam}{int}{pr}
\olsection{Fungsi kang \usetoken{S}{lambda definable} Katutup tumrap Rekursi Primitif}

Nalika tekan rekursi primitif, pungkasane ana
pakaryan kang kudu ditindakake. Kita kudu
lumaku kanthi sawetara tahap. Kaya sadurunge,
upamana wis ana suku $G'$ lan $H'$ kang
saben-saben !!{lambda define} fungsi $g$
lan~$h$. Kita arep nemokake suku $F$ kang
!!{lambda define}s fungsi~$f$ kang ditetepake
dening
\begin{align*}
f(0, \vec z) & = g(\vec z) \\
f(x+1, \vec z) & = h(x, f(x,\vec z), \vec z).
\end{align*}
Dadi, umume, yen ana suku lambda $G'$
lan~$H'$, cukup nemokake suku~$F$ saengga
\begin{align*}
F(\num{0}, \vec z) & \equiv G'(\vec z) \\
F(\overline{n+1}, \vec z) & \equiv H'(\num{n}, F(\num{n}, \vec z), \vec z)
\end{align*}
kanggo saben wilangan asli~$n$; amarga
$G'$ lan~$H'$ !!{lambda define} $g$
lan~$h$, yen numeral $\num{\vec m}$
diselehake kanggo $\vec z$,
$F(\num{n+1}, \num{\vec m})$ bakal
direduksi nganti wujud normal kang
menehi asil kang bener.

Nanging kanggo iki, cukup nemokake
suku~$F$ kang nyukupi
\begin{align*}
F(\num 0) & \equiv G \\
F(\overline {n+1}) & \equiv H(\num{n},F(\num{n}))
\intertext{kanggo saben wilangan asli $n$, kanthi}
G & = \lambd[\vec z][G'(\vec z)] \text{ lan}\\
H(u,v) & = \lambd[\vec z][H'(u,v(\vec z),\vec z)].
\end{align*}
Tegese, kanthi akal-akalan notasi lambda,
kita ora perlu maneh mikirake paramèter
tambahan $\vec z$: paramèter iku dilebokake
ing notasi abstraksi lambda.

Sadurunge netepake suku $F$, kita perlu
cara kanggo nangani pasangan terurut.
Lema sabanjure nyedhiyakake cara iku.

\begin{lem}
Ana suku lambda $D$ saengga kanggo saben
pasangan suku lambda $M$ lan~$N$,
$D(M,N)(\num{0}) \red M$ lan
$D(M,N)(\num{1}) \red N$.
\end{lem}

\begin{proof}
Dhisik, tetepna suku lambda $K$ kanthi
\[
K(y) = \lambd[x][y].
\]
Tegese, $K$ iku suku
$\lambd[y][\lambd[x][y]]$.
Kanthi cara ndeleng liyane, kanggo saben
$M$, $K(M)$ iku fungsi konstan kang
ngasilake~$M$ kanggo input apa wae.

Saiki tetepna $D(x,y,z)$ kanthi
$D(x,y,z) = z (K(y))x$. Banjur kita duwe
\begin{align*}
D(M,N,\num 0) & \red \num 0 (K(N)) M \red M \text{ lan}\\
D(M,N,\num 1) & \red \num 1 (K(N)) M \red K(N) M \red N,
\end{align*}
kaya kang dikarepake.
\end{proof}

Gagasane yaiku $D(M,N)$ makili pasangan
$\tuple{M, N}$; yen $P$ dianggep makili
pasangan kaya mangkene, $P(\num 0)$
lan $P(\num 1)$ makili proyeksi kiwa lan
tengen, yaiku $(P)_0$ lan $(P)_1$.
Kita bakal nggunakake notasi pungkasan iki.

\begin{lem}
Fungsi-fungsi kang !!{lambda definable}
katutup tumrap rekursi primitif.
\end{lem}

\begin{proof}
Kita kudu nuduhake yen, yen ana
sembarang suku $G$ lan $H$, kita bisa
nemokake suku $F$ saengga
\begin{align*}
F(\num{0}) & \equiv G \\
F(\num{n+1}) & \equiv H(\num{n}, F(\num{n}))
\end{align*}
kanggo saben wilangan asli~$n$.
Gagasane, kanthi kasar, yaiku ngitung
urutan \emph{pasangan}
\[
\tuple{\num 0, F(\num 0)}, \tuple{\num 1, F(\num 1)}, \dots,
\]
kanthi nggunakake numeral minangka iterator.
Gatekna yen pasangan kapisan mung
$\tuple{\num 0, G}$. Yen wis ana pasangan
$\tuple{\num{n}, F(\num{n})}$, pasangan
sabanjure $\tuple{\num{n+1}, F(\num{n+1})}$
kudu ekuivalen karo
$\tuple{\num{n+1}, H(\num{n}, F(\num{n}))}$.
Kita bakal ngrancang suku lambda~$T$
kang nindakake transisi saklangkah iki.

Rinciane kaya mangkene. Tetepna $T(u)$
kanthi
\[
T(u) = \tuple{\fn{Succ}((u)_0), H((u)_0,(u)_1)}.
\]
Saiki gampang dipriksa yen kanggo saben
wilangan~$n$,
\[
T(\tuple{\num{n}, M}) \red \tuple{\num{n+1}, H(\num{n}, M)}.
\]
Kaya kang disaranake ing ndhuwur, yen
wis ana $G$ lan $H$, tetepna~$F(u)$
kanthi
\[
F(u) = (u(T,\tuple{\num 0, G}))_1.
\]
Tegese, kanggo input~$\num{n}$,
$F$ ngetrapake $T$ ping $n$ marang
$\tuple{\num 0, G}$, banjur ngasilake
komponen kapindho. Kanggo wiwitane,
kita duwe
\begin{enumerate}
\item $\num{0} (T, \tuple{\num 0, G}) \equiv \tuple{\num{0}, G}$
\item $F(\num{0}) \equiv G$
\end{enumerate}
Kanthi induksi ing~$n$, kita bisa
nuduhake yen kanggo saben wilangan asli
ukara ing ngisor iki bener:
\begin{enumerate}
\item $\num{n+1}(T, \tuple{\num{0}, G}) \equiv \tuple{\num{n+1}, F(\num{n+1})}$
\item $F(\num{n+1}) \equiv H(\num{n}, F(\num{n}))$
\end{enumerate}
Kanggo klausa kapindho, kita duwe
\begin{align*}
F(\num{n+1}) & \red (\num{n+1}(T, \tuple{\num 0, G}))_1 \\
& \equiv (T(\num{n} (T, \tuple{\num 0, G})))_1 \\
& \equiv (T(\tuple{\num{n}, F(\num{n})}))_1 \\
& \equiv (\tuple{\num{n+1}, H(\num{n}, F(\num{n}))})_1 \\
& \equiv H(\num{n}, F(\num{n})).
\end{align*}
Ing larik sadurunge larik pungkasan,
kita nggunakake hipotesis induksi.
Kanggo klausa kapisan, kita duwe
\begin{align*}
\num{n+1} (T, \tuple{\num 0, G}) &
\equiv T(\num{n} (T, \tuple{\num 0, G})) \\
& \equiv T( \tuple{\num{n}, F(\num{n})}) \\
& \equiv \tuple{\num{n+1}, H(\num{n}, F(\num{n}))} \\
& \equiv \tuple{\num{n+1}, F(\num{n+1})}.
\end{align*}
Ing larik pungkasan, kita nggunakake
klausa kapindho. Dadi kita wis nuduhake
$F(\num 0) \equiv G$ lan, kanggo saben
$n$, $F(\num {n+1}) \equiv H(\num
n, F(\num{n}))$, persis kaya kang
dibutuhake.
\end{proof}

\end{document}
